\documentclass[11pt]{article}
\usepackage{guide}
\guideheader{Deep dive 2: the eBPF tracer}
\begin{document}

\guidetitle{Watching the Kernel at Work}{Deep dive 2: the eBPF tracer, the collector and what tracing costs}{For Rohit. Read the overview first.}

\section{Your segment in one page}

\pipeline{k,c}

To explain why an operation was slow we need evidence from inside the kernel: when did the thread stop running, why, when was it woken, what did the disk do meanwhile? Your segment collects that evidence. It has two halves:

\begin{itemize}
\item \textbf{The eBPF programs} run \emph{inside} the kernel. Each one is attached to a point of interest, and writes one small record every time that point is reached.
\item \textbf{The collector} runs as a normal program. It loads the eBPF programs, starts the workload, receives the records and writes them to a file.
\end{itemize}

\begin{center}
\begin{tabularx}{\textwidth}{@{}lY@{}}
\toprule
Goes in & A command to run (the workload) \\
Comes out & \f{events.bin}: every event as a 48-byte record. \f{stats.json}: what was attached, how many events were produced, received and lost, and how much CPU the collector used \\
Kernel side & \f{src/ebpf/stalltrace.bpf.c} (about 490 lines), \f{src/ebpf/stalltrace.h} \\
User side & \f{src/collector/stalltrace.c} (about 770 lines) \\
Build and cost & \f{Makefile}, \f{scripts/run\_overhead.sh}, \f{src/classifier/overhead.py} \\
\bottomrule
\end{tabularx}
\end{center}

\section{Concepts you need}

\subsection{eBPF and the verifier}
eBPF lets a program run small pieces of code inside the Linux kernel without changing or rebuilding the kernel. Because a bug in the kernel can crash the whole machine, every eBPF program is checked before loading by the \emph{verifier}. It rejects anything it cannot prove safe: loops without a bound, reading memory through an unchecked pointer, and so on. Getting past the verifier was the hardest part of this segment (section~5).

\subsection{Where a program can attach}
\begin{itemize}
\item \textbf{Tracepoints} are stable hooks the kernel developers placed at important events, such as ``the scheduler is switching threads''. We use them for system calls, the scheduler, the block layer and page faults.
\item \textbf{Uprobes} attach to a function in a normal program. We use two, on the workload's \f{op\_begin} and \f{op\_end} markers, so we know where each operation starts and ends.
\end{itemize}

\subsection{CO-RE: compile once, run everywhere}
Kernel data structures change between versions. A field may move or disappear. \textbf{CO-RE} solves this: the compiled eBPF program records \emph{which field} it wants, and the loader looks up where that field is in the running kernel, using type information the kernel publishes (called BTF, at \f{/sys/kernel/btf/vmlinux}). The library that does this is \textbf{libbpf}. This is why our tracer needs no kernel headers and can check at load time whether a field such as \f{in\_iowait} exists.

\subsection{Maps and the ring buffer}
eBPF programs cannot print or write files. They share data through \textbf{maps}. We use three kinds:
\begin{itemize}
\item a \textbf{hash map} listing which processes we are tracing;
\item a \textbf{ring buffer}, a fixed-size queue shared with the collector. The kernel side \emph{reserves} space, fills in the record and \emph{submits} it; the collector reads from the other end;
\item \textbf{per-CPU counters} of events produced and dropped.
\end{itemize}

If the collector reads more slowly than the kernel writes, the ring buffer fills up and new events are \emph{dropped}. A trace with silently missing events gives wrong answers, so we count every drop and report the count for every run.

\subsection{\f{in\_iowait}: the kernel's own hint}
When a thread goes to sleep waiting for disk I/O, the kernel sets a one-bit flag on it called \f{in\_iowait}. Reading that flag at the moment the thread leaves the CPU tells us whether the sleep is for the disk or for something else, such as a timer. It turned out to be the single most useful piece of evidence in the whole project.

\subsection{PID namespaces}
A process can have different ID numbers depending on where you look from. Containers, and also WSL, run processes in a \emph{PID namespace}: the workload thinks it is process 1592, while the kernel as a whole knows it as 3614. eBPF helpers return the kernel-wide number. We must translate, or we trace the wrong process (section~5).

\section{Code walkthrough}

\subsection{One event format for everything}
\src{src/ebpf/stalltrace.h}
\begin{lstlisting}[language=C]
struct st_event {
	__u64 ts_ns;          /* when (monotonic clock, same as the workload) */
	__u32 pid, tid, cpu;  /* who and where */
	__u16 type;           /* which of the 11 event types */
	__u16 flags;          /* preempted? in_iowait? read or write? ... */
	__u64 a, b;           /* meaning depends on the type */
	__u32 dev, c;
};                        /* exactly 48 bytes */
\end{lstlisting}
Every event, whatever its kind, is the same small fixed-size record. Small records mean fewer drops. The meaning of \f{a}, \f{b} and \f{c} depends on the type: for a system call \f{a} is the call number; for a disk request \f{a} is the sector and \f{b} the size in bytes; for a marker \f{a} is the operation ID.

\subsection{Reserving space and counting drops}
\src{src/ebpf/stalltrace.bpf.c, reserve(), lines 176--198 (shortened)}
\begin{lstlisting}[language=C]
static __always_inline struct st_event *reserve(__u32 type)
{
	struct st_event *e;

	bump(&produced, type);                         /* count every attempt */
	e = bpf_ringbuf_reserve(&events, sizeof(*e), 0);
	if (!e) {
		bump(&dropped, type);                      /* buffer full: count the loss */
		return NULL;
	}
	e->ts_ns = bpf_ktime_get_ns();
	e->cpu   = bpf_get_smp_processor_id();
	e->type  = type;
	return e;
}
\end{lstlisting}
Every handler calls this first. It is where dropped events are detected: if the ring buffer has no room, the reservation fails and we count it, separately for each event type.

\subsection{The scheduler hook: the most important one}
\src{handle\_sched\_switch(), lines 303--352 (shortened)}
\begin{lstlisting}[language=C]
SEC("tp_btf/sched_switch")
int BPF_PROG(handle_sched_switch, bool preempt, struct task_struct *prev,
	     struct task_struct *next, unsigned int prev_state)
{
	/* ... skip unless prev or next belongs to a traced process ... */
	e = reserve(ST_EV_SCHED_SWITCH);
	e->a = prev_state;                             /* runnable or sleeping? */
	if (preempt)
		e->flags |= STF_PREEMPT;                   /* kicked off the CPU */
	if (BPF_CORE_READ(prev, pid) == 0)
		e->flags |= STF_PREV_IDLE;                 /* the CPU was idle before */
	if (bpf_core_field_exists(prev->in_iowait)) {
		e->flags |= STF_IOWAIT_KNOWN;
		if (BPF_CORE_READ_BITFIELD_PROBED(prev, in_iowait))
			e->flags |= STF_IOWAIT;                /* sleeping for the disk */
	}
	bpf_ringbuf_submit(e, 0);
	return 0;
}
\end{lstlisting}
This runs every time a CPU switches from one thread (\f{prev}) to another (\f{next}). For the thread leaving, we record three facts that decide everything later:
\begin{itemize}
\item \textbf{Was it preempted?} Then it still wants to run and is waiting for a CPU: scheduler wait.
\item \textbf{Did it go to sleep with \f{in\_iowait} set?} Then it is waiting for the disk: I/O wait.
\item \textbf{Did it go to sleep without it?} Then it is blocked on something else.
\end{itemize}
\f{bpf\_core\_field\_exists} is CO-RE at work: if a kernel does not have the field, we say so (\f{STF\_IOWAIT\_KNOWN} stays off) instead of guessing.

\subsection{The other hooks}
\begin{center}
\begin{tabularx}{\textwidth}{@{}lY@{}}
\toprule
Hook & What it records \\
\midrule
\f{raw\_syscalls/sys\_enter}, \f{sys\_exit} & Start and end of every system call by the traced process \\
\f{sched\_wakeup} & A sleeping traced thread became ready to run \\
\f{block\_bio\_queue} & A thread submitted a disk request (this is how a request is tied to the thread that caused it) \\
\f{block\_rq\_issue}, \f{block\_rq\_complete} & A request was sent to the device, and finished. Recorded for \emph{all} processes on the test disk, so another program's I/O stays visible \\
\f{exceptions/page\_fault\_user} & A page fault (needed for the \f{mmap} workloads) \\
uprobes on \f{op\_begin}, \f{op\_end} & Start and end of each operation, with its ID \\
\bottomrule
\end{tabularx}
\end{center}

\subsection{Knowing which process to trace}
\src{try\_register(), lines 111--127}
\begin{lstlisting}[language=C]
leader = BPF_CORE_READ(t, group_leader);
if (task_ns_pid(leader, &ns_tgid, &inum) || inum != pidns_inum)
	return NULL;                                   /* not in our namespace */
if (!bpf_map_lookup_elem(&pending_tgids, &ns_tgid))
	return NULL;                                   /* not the one we wait for */
bpf_map_update_elem(&traced_tgids, &host_tgid, &ns_tgid, BPF_ANY);
\end{lstlisting}
The collector only knows the workload's process ID as seen inside its namespace. It puts that number in the \f{pending\_tgids} map. The first time the kernel side sees a process whose \emph{namespace} ID matches, it records the pair (kernel-wide ID, namespace ID) in \f{traced\_tgids}. From then on, lookups use the kernel-wide ID and events are written with the namespace ID, so they match what the workload itself reports.

\subsection{The collector}
\src{src/collector/stalltrace.c, main(), simplified}
\begin{lstlisting}[language=C]
skel = stalltrace_bpf__open();                     /* parse the eBPF object */
bpf_map__set_max_entries(skel->maps.events, opt.rb_bytes);   /* ring size (-b) */
stalltrace_bpf__load(skel);                        /* verifier runs here */
/* attach tracepoints, then uprobes on op_begin / op_end in the workload */
rb = ring_buffer__new(bpf_map__fd(skel->maps.events), handle_event, &out, NULL);

target = fork();                                   /* child waits at a "gate" */
bpf_map_update_elem(pending_tgids, &target, ...);  /* tell the kernel side */
write(gate[1], "g", 1);                            /* open the gate: child execs */

for (;;) {
	ring_buffer__poll(rb, 100);                    /* handle_event() per record */
	if (workload exited and linger time passed) break;
}
\end{lstlisting}
Two details are worth noticing. First, the \emph{gate}: the child process is created but held back until tracing is fully set up, so no early event is missed. Second, \f{handle\_event()} does almost nothing: it copies the 48 bytes straight to the file. An earlier version formatted each event as text, and that alone made the collector too slow (section~5).

\f{stalltrace.skel.h}, used above, is generated by \f{bpftool} from the compiled eBPF object during the build. The \f{Makefile} does it in four steps: dump kernel types to \f{vmlinux.h}, compile the eBPF C to an object with \f{clang -target bpf}, generate the skeleton, compile the collector.

\section{The paper behind it: No Two Snowflakes Are Alike (eBPF '25)}

\textbf{What it does.} Machado and colleagues implement the same three tracing tools with five eBPF libraries (bpftrace, BCC, libbpf, ebpf-go and Aya) and measure, under heavy storage workloads, how much each slows the application, how much CPU and memory it uses, and how many events it loses.

\textbf{What they found.} The libraries differ widely. In their read/write tracer test, which generates about 33 million events, libbpf captured all of them, Aya 98\%, BCC 47\%, ebpf-go 37\% and bpftrace only 14\%. When each event carried a large payload (about 4\,KiB), every library lost more than 99\%.

\textbf{What we took from it.}
\begin{itemize}
\item \emph{Use libbpf.} It lost the fewest events and used the least CPU and memory.
\item \emph{Keep events small.} Hence our fixed 48-byte record.
\item \emph{Measure loss, do not assume it is zero.} Hence the produced and dropped counters, reported for every run.
\item \emph{Their settings as our starting point:} a 256\,KiB ring buffer polled every 100\,ms, on Ubuntu 24.04.
\end{itemize}

\textbf{Where we differ.} They study the tracer itself. We use a tracer to diagnose applications, and we report its fidelity as part of our results.

\section{Problems we hit and how we fixed them}

\begin{itemize}
\item \textbf{The verifier rejected the block-layer handler.} Error: ``dereference of modified ctx ptr''. We were reading a text field of the tracepoint through an adjusted pointer. Fix: copy the field to the stack first with \f{bpf\_probe\_read\_kernel}.
\item \textbf{The verifier rejected the scheduler handler.} Error: ``bitwise operator |= on pointer''. The compiler had combined two map lookups by OR-ing their pointers. Fix: copy the values out of the map into plain variables immediately (the comment at the top of \f{handle\_sched\_switch} explains this).
\item \textbf{Only seven events were captured.} The first working version traced almost nothing in WSL because of the PID namespace mismatch. Fix: the registration scheme above.
\item \textbf{Waking an idle CPU looks like waiting.} In a virtual machine, a woken thread starts running 30--40 microseconds later because the CPU must leave its idle state. Fix: the \f{STF\_PREV\_IDLE} flag marks these cases so the analysis does not count them as CPU contention.
\item \textbf{The collector could not keep up.} Writing each event as a line of JSON text cost so much CPU that up to half the events of fast workloads were dropped. Fix: write the raw 48-byte records (\f{events.bin}); text output is still available with \f{-J}.
\end{itemize}

\section{Results that come from this segment}

\textbf{Fidelity in the final experiment:} 766{,}242 events, 0.54\% dropped. Scores do not change when runs with drops are removed.

\textbf{Overhead} (tracing on versus off, median of five runs):

\begin{center}
\begin{tabular}{@{}lcrrr@{}}
\toprule
Workload & Threads & Slowdown & Added CPU per op & Dropped \\
\midrule
\f{sleep} & 1 & 4.3\% & 139\,\textmu s & 0\% \\
\f{sync\_write} & 1 & 1.1\% & 195\,\textmu s & 0\% \\
\f{direct\_read} & 4 & 13.3\% & 81\,\textmu s & 0\% \\
\f{warm\_read} & 4 & 61.1\% & 21\,\textmu s & 17.6\% \\
\f{buffered\_write} & 1 & 91.5\% & 27\,\textmu s & 0.5\% \\
\midrule
\f{warm\_read}, 4\,MiB ring buffer & 4 & 73.6\% & 32\,\textmu s & 0\% \\
\f{warm\_read}, no markers & 4 & 33.0\% & 7\,\textmu s & 0.6\% \\
\bottomrule
\end{tabular}
\end{center}

How to read this:
\begin{itemize}
\item The cost is roughly a \emph{fixed amount per operation}. For an operation that lasts milliseconds it is a few percent. For one that lasts 10 microseconds it is several times the operation itself.
\item Most of the fixed cost is the two marker uprobes, about 7--10 microseconds each in WSL. Removing them cuts the added CPU for \f{warm\_read} from 21 to 7\,\textmu s.
\item Drops are a buffer-size problem: a 4\,MiB ring buffer removes them completely.
\item Our budget, declared before measuring, was at most 10\% throughput loss and under 1\% drops. We meet it for slow operations and miss it for fast ones, and the report says so.
\item These measurements are noisy: two identical untraced disk-read runs differed by 15--41\%.
\end{itemize}

\section{Try it yourself}

\textbf{Without administrator rights}, look at what the tracer recorded. Download \f{raw-wsl.tar.gz} from the GitHub release into the project folder, then:

\begin{lstlisting}[language=bash]
mkdir -p raw && tar xzf raw-wsl.tar.gz -C raw
R=raw/final-wsl/direct_read-t4-io64k-seq-d2ms-r3
python3 -c "import json; s=json.load(open('$R/stats.json')); print(s['programs']); print(s['totals'])"
grep -m1 -A2 '"iowait":1' $R/events.jsonl
\end{lstlisting}

You should see twelve programs marked \f{attached} and totals of 2{,}038 events with 0 dropped. The last command prints three events that together are one disk wait: a \f{sched\_switch} where a worker thread leaves the CPU with \f{iowait} set to 1, a \f{sched\_wakeup} for the same thread when its I/O completes, and a \f{sched\_switch} that puts it back on the CPU with \f{prev\_idle} set to 1.

\textbf{With administrator rights} (Linux or WSL):

\begin{lstlisting}[language=bash]
sudo scripts/setup_vm.sh && make
sudo ./build/stalltrace -o /tmp/t1 -J -- ./build/stallbench --workload sleep \
    --run-id t1 --out /tmp/t1.gt.jsonl --ops 5
cat /tmp/t1/stats.json
\end{lstlisting}

\f{./build/stalltrace -h} lists the options: \f{-b} sets the ring buffer size, and \f{-S}, \f{-C}, \f{-K}, \f{-F}, \f{-M} switch off system-call, scheduler, block, page-fault and marker events.

\section{Limits and open questions}

\begin{itemize}
\item \textbf{The markers are expensive.} A lighter way to mark operation boundaries (USDT probes, or boundaries taken from kernel events) would widen the range of workloads the tool suits.
\item \textbf{Page-fault tracing is x86 only} in the current code.
\item \textbf{Needs a recent kernel and root.} The scheduler hook uses the argument layout of kernel 5.18 and later, and loading eBPF needs administrator rights.
\item \textbf{Tested on one kernel.} WSL2's 6.6 kernel. Whether everything attaches on a stock Ubuntu 6.8 kernel is part of the pending robustness check; \f{scripts/check\_capabilities.sh} reports what is available.
\end{itemize}

\end{document}
